\begin{textsample}{POS Dim 3 – persona\_grok – Score 18.00 – t850\_grok.txt}  \label{ex:f3_pos_033}
Since 2007, the smartphone \textbf{industry} has produced 7.1 \textbf{billion} devices, each relying on over 60 different \textbf{elements}.
This massive \textbf{production} has driven extensive mining and processing activities, resulting in profound environmental and social impacts around the world.
The consequences are stark.
In 2014 alone, the world generated 3 million metric \textbf{tons} of electronic \textbf{waste} from smartphones and similar devices, but only 16 % of that e-waste was properly recycled.
Most smartphone \textbf{models} don't feature easily replaceable batteries, which shortens their usable lifespan to about two years in the US, exacerbating the \textbf{waste} \textbf{problem}. \textbf{Energy} demands are equally alarming : manufacturing these \textbf{billions} of smartphones has consumed 968 terawatt-hours of \textbf{energy}, an \textbf{amount} \textbf{equivalent} to India 's entire annual \textbf{electricity} \textbf{supply}.
Inefficient \textbf{recycling} practices compound these issues, as seen in high-profile cases like the Samsung Galaxy Note 7 recall, where faulty devices highlighted the flaws in current \textbf{production} and \textbf{waste} management \textbf{systems}.
It 's \textbf{time} for a fundamental \textbf{shift} toward a circular \textbf{economy} \textbf{model} that prioritizes sustainability.
We urge supporters to join petitions calling on Samsung to improve repairability and \textbf{recycling} in their \textbf{products}, helping to reduce the environmental \textbf{footprint} of our tech-driven world.

% matched lemmas: amount, billion, economy, electricity, element, energy, equivalent, footprint, industry, model, problem, product, production, recycling, shift, supply, system, time, ton, waste
\end{textsample}
